%%%%%%%%%%%%%%%%%%%%%%%%%%%%%%%%%%%%%%%%%%%%%%%%%%%%%%%%%%%%%%%%%
\chapter{SİSTEM MİMARİSİ}\label{ch:sistem_mimarisi}
%%%%%%%%%%%%%%%%%%%%%%%%%%%%%%%%%%%%%%%%%%%%%%%%%%%%%%%%%%%%%%%%%
\vspace*{-1cm}

Bu proje, modern finansal takip ve yapay zekâ tabanlı bütçe optimizasyonu sunan bir web uygulamasıdır. Bu bölümde, sistemin genel yapısı, kullanılan teknolojiler ve veri modelleri detaylandırılmıştır.

\section{Genel Sistem Mimarisi}

Sistem, bir monorepo yapısında tasarlanmış olup üç ana bileşenden oluşmaktadır:
\begin{enumerate}
    \item \textbf{İstemci (Frontend):} React (Vite) tabanlı kullanıcı arayüzü.
    \item \textbf{Sunucu (Backend):} Node.js ve Express.js tabanlı API servisleri.
    \item \textbf{Yapay Zekâ Servisi:} Python tabanlı bütçe optimizasyon ve tahmin modelleri.
\end{enumerate}

Şekil~\ref{fig:mimari_diyagram} sistemin genel çalışma prensibini göstermektedir. Kullanıcılar web arayüzü üzerinden harcamalarını girerken, sunucu bu verileri PostgreSQL veritabanına kaydeder ve gerektiğinde Python tabanlı ML modellerini veya Node.js tarafındaki LSTM ağını tetikler.

\begin{figure}[h!]
  \centering
  \includegraphics[width=0.8\textwidth]{aura_sc1}
  \caption{Aura Finance Genel Sistem Mimarisi.}
  \label{fig:mimari_diyagram}
\end{figure}

\section{Veritabanı Şeması}

Sistemin kalbinde PostgreSQL veritabanı yer almaktadır. Veri modelleri, Prisma ORM kullanılarak tanımlanmış ve yönetilmiştir. Temel tablolar ve ilişkileri aşağıda açıklanmıştır:

\subsection{Kullanıcı (User) Modeli}
Kullanıcıların temel bilgilerini (ad, e-posta, şifre karması) ve aylık toplam gelir (income) bilgisini saklar.

\subsection{Harcama (Expense) Modeli}
Her bir harcamanın miktarını, kategorisini, tarihini ve açıklamasını içerir. Ayrıca harcamanın kişisel mi yoksa bir gruba mı ait olduğunu belirten \texttt{groupId} alanına sahiptir.

\subsection{Bütçe (Budget) Modeli}
Kullanıcıların her bir kategori için belirlediği limitleri saklar. Yapay zekâ tarafından önerilen bütçeler de bu tabloya güncellenebilmektedir.

\subsection{Grup ve Grup Üyeliği (Group \& GroupMember)}
Paylaşımlı harcamalar için tasarlanmıştır. Kullanıcılar gruplar oluşturabilir ve harcamaları grup üyeleri arasında otomatik olarak paylaştırabilirler.

\section{Teknoloji Yığını}

\subsection{İstemci Tarafı Teknolojileri}
Aura Finance istemcisi, yüksek performans ve tip güvenliği için \textbf{TypeScript} ile yazılmış bir \textbf{React} uygulamasıdır. \textbf{Tailwind CSS} ve \textbf{Shadcn UI} bileşenleri kullanılarak modern ve duyarlı (responsive) bir arayüz tasarlanmıştır. Veri yönetimi için \textbf{React Query} kullanılarak sunucu ile senkronizasyon optimize edilmiştir.

\subsection{Sunucu Tarafı Teknolojileri}
Arka uçta \textbf{Node.js} ve \textbf{Express.js} kullanılmıştır. Kimlik doğrulama işlemleri için \textbf{JWT (JSON Web Token)} ve \textbf{bcryptjs} kütüphaneleri tercih edilmiştir. Veritabanı işlemleri için \textbf{PostgreSQL} tercih edilmiş, şema yönetimi ve sorgular için \textbf{Prisma} ve \textbf{pg} kütüphaneleri kullanılmıştır.

\subsection{Yapay Zekâ Teknolojileri}
\begin{itemize}
    \item \textbf{TensorFlow.js:} Node.js ortamında zaman serisi tahmini (LSTM) yapmak için kullanılmıştır.
    \item \textbf{XGBoost (Python):} Karmaşık bütçe optimizasyon algoritmaları için Multi-Output Regressor modeli olarak uygulanmıştır.
    \item \textbf{Tesseract.js:} İstemci tarafında OCR (Optik Karakter Tanıma) ile makbuz okuma özelliği sağlanmıştır.
\end{itemize}
